%-----------------------------------------------------------------------------------------------
\makeatletter
\immediate\write18{datelog > \jobname.info} % site script for $(date '+%Y-%m-%d %Hh%Mm%Ss')
\makeatother
%-----------------------------------------------------------------------------------------------
%-----------------------------------------------------------------------------------------------
\usetheme{Copenhagen}
\usepackage{beamercolorthemeCNThermSci2}
\usefonttheme{serif}
%-----------------------------------------------------------------------------------------------

%-----------------------------------------------------------------------------------------------
%-----------------------------------------------------------------------------------------------
\usetheme{Copenhagen}
\usepackage{beamercolorthemeUTF2}
\usefonttheme{serif}
%-----------------------------------------------------------------------------------------------
\usepackage[utf8]{inputenc}
\usepackage[greek,french,english,brazil]{babel} % last becomes the active one
\usepackage{pslatex}
\usepackage{amssymb,amsmath}
\usepackage{soul}
\usepackage[squaren,Gray]{SIunits}
\usepackage[nice]{nicefrac}
\usepackage{tikz}
\usepackage{xspace}
%-----------------------------------------------------------------------------------------------


%-----------------------------------------------------------------------------------------------
%-----------------------------------------------------------------------------------------------
% Mathematical
%-----------------------------------------------------------------------------------------------
\newcommand{\vet}[1]{\underline{{#1}}}
\newcommand{\mat}[1]{\underline{\underline{{#1}}}}
\newcommand{\cub}[1]{\underline{\underline{\underline{{#1}}}}}
\newcommand{\eqdef}{{\ensuremath\stackrel{\text{\tiny def}}{=}}}
%-----------------------------------------------------------------------------------------------
% Presentation
%-----------------------------------------------------------------------------------------------
\newcommand{\BkgImgH}[1]{% Places an image centered on the slide background filling the height
    \usebackgroundtemplate{\parbox{\paperwidth}{%
        \vspace*{1sp}\centering\includegraphics[height=\paperheight]{{#1}}
}}}
\newcommand{\BkgImgW}[1]{% Places an image centered on the slide background filling the width
    \usebackgroundtemplate{\parbox{\paperwidth}{%
        \vspace*{1sp}\centering\includegraphics[width=\paperwidth]{{#1}}
}}}
\newcommand{\ArtEndH}[3]{% Transitions to plain image (last) slide: #1:prefix #2,#3:extensions
    \BkgImgH{root/../art/#1.#2}
    \frame[plain]<handout:0>{%
        \transdissolve\vspace*{80mm}\color{white}\bf\input{root/../art/#1.#3}
}}
\newcommand{\ArtEndW}[3]{% Transitions to plain image (last) slide: #1:prefix #2,#3:extensions
    \BkgImgW{root/../art/#1.#2}
    \frame[plain]<handout:0>{%
        \transdissolve\vspace*{80mm}\color{white}\bf\input{root/../art/#1.#3}
}}
\newcommand{\ImgColW}[3]{% Inserts a full-width image in a column
    \includegraphics[width=\columnwidth]{root/../art/#1.#2}\\[-0.5\baselineskip]
    \parbox{\columnwidth}{\tiny\hfill\scalebox{0.65}{\input{root/../art/#1.#3}}}
}
\newcommand{\txtpic}[1]{%
    \fcolorbox{lightgray}{white!90!black}{{#1}} 
}
%-----------------------------------------------------------------------------------------------


%-----------------------------------------------------------------------------------------------
\title{C.02.01 -- Ciclo Otto Ar-Combustível de Tempo Finito de Combustão}
\subtitle{FTAF -- Finite Time Air-Fuel Otto Engine Model}
\author{Prof.~C.~Naaktgeboren, PhD}
\date{{\scriptsize\tt%
    \includegraphics[height=6.0mm]{root/00-res/cc/by-nc-nd-88x31.pdf}\\[\smallskipamount]
    https://github.com/CNThermSci/ApplThermSci\\
    Compiled on \input{\jobname.info}
}}
%-----------------------------------------------------------------------------------------------
\begin{document}
%-----------------------------------------------------------------------------------------------
\logo{%
    \parbox{158mm}{% There's a 1mm gap on each side of the 160mm x 90mm slide logo line
        \mode<beamer>{
            \includegraphics[height=6.0mm]{root/00-res/UTFPR/UTFPR-logo-D.pdf}\hfill%
            \includegraphics[height=9.0mm]{root/00-res/logo/CNThermSci-logo-A.pdf}%
        }
        \mode<handout>{
            \includegraphics[height=6.0mm]{root/00-res/UTFPR/UTFPR-logo-W.pdf}\hfill%
            \includegraphics[height=9.0mm]{root/00-res/logo/CNThermSci-logo-W.pdf}%
        }
    }
} % The (delineated, alpha), or washed-out logos
%-----------------------------------------------------------------------------------------------
\frame{\titlepage}
%-----------------------------------------------------------------------------------------------

%------------------------------------------------------------------------------------------------
\section{Apresentação do FTAF}
%-----------------------------------------------------------------------------------------------

    % !j 96 -i8
    %-------------------------------------------------------------------------------------------
    \begin{frame}{Ciclo Otto ar-combustível de tempo finito---FTAF}\vspace*{-2em}
        \begin{itemize}
            \item<1->  Modelo do livro-texto (tópicos de leitura) adiciona combustão ao
                \alert{Ciclo Otto ideal};
            \begin{itemize}
                \item<2->  Permite variação de \alert{combustíveis};
                \item<3->  Porém, desde que sejam \alert{carbonados}: norm.~em $C$; excluindo
                    \ce{H2}, \ce{H4N2} puros, p.~ex.;
                \item<3->  Ênfase nas \alert{propriedades} $\bar{c}_{p,v}(T)$, $k(T)$,
                    $\bar{u}(T)$, etc.~das misturas;
                \item<4->  Incorpora \alert{combustão} e \alert{equilíbrio químico};
                \item<5->  \alert{Não emprega} o calor liberado na \alert{combustão}!
            \end{itemize}
            \item<6->  Modelo \alert{ar-combustível de tempo finito, FTAF}:
            \begin{itemize}
                \item<7->  Adiciona \alert{combustão}, mantendo as demais características do
                    \alert{FTHA};
                \item<8->  Obtém tanto as \alert{propriedades} quanto o \alert{calor liberado}
                    pelas \alert{reações}!
                \item<9->  Desenvolvido em um \alert{TCC} defendido em \alert{2018} (citação nos
                    tópicos de leitura);
            \end{itemize}
        \end{itemize}
    \end{frame}
    %-------------------------------------------------------------------------------------------

%------------------------------------------------------------------------------------------------
\subsection{Como Extensão do FTAH}
%-----------------------------------------------------------------------------------------------

    % !j 96 -i8
    %-------------------------------------------------------------------------------------------
    \begin{frame}{Ciclo Otto padrão a ar de tempo finito de adição de calor---FTHA}\vspace*{-2em}
        \begin{itemize}
            \item<1->  Modela combustão (adição de calor) de forma \alert{não instantânea}:
            \begin{itemize}
                \item<2->  Interações \alert{simultâneas} de \alert{calor} e \alert{trabalho};
                \item<3->  Tempos de motor \alert{discretizados} em \alert{sub-processos};
                \item<4->  Elemento computacional: sub-processo \alert{localmente politrópico};
                \item<5->  \alert{Remoção} de calor permanece \alert{isocórica} (instantânea).
            \end{itemize}
            \item<6->  Mantém-se como modelo \alert{padrão a ar}:
            \begin{itemize}
                \item<7->  Transferência de calor para bloco inclui \alert{irreversibilidades};
                \item<8->  Perdas de bombeamento envolvem \alert{sistema e ciclo abertos}.
            \end{itemize}
            \item<9->  Mantém-se como modelo de \alert{substância pura}:
            \begin{itemize}
                \item<10-> Evita \alert{combustão e equilíbrio químico};
                \item<11-> Evita modelagem termodinâmica de \alert{misturas reativas}.
            \end{itemize}
        \end{itemize}
    \end{frame}
    %-------------------------------------------------------------------------------------------

    % !j 96 -i8
    %-------------------------------------------------------------------------------------------
    \begin{frame}{Ciclo Otto padrão a ar de tempo finito de adição de calor---FTHA}\vspace*{-2em}
        \begin{itemize}
            \item<1->  Inclui todos os parâmetros do \alert{ciclo Otto ideal}:
            \begin{itemize}
                \item<2->  \alert{Razão de compressão} do motor;
                \item<3->  \alert{Calores específicos} do fluido de trabalho.
            \end{itemize}
            \item<4->  Inclui parâmetros \alert{construtivos} do \alert{motor}:
            \begin{itemize}
                \item<5->  Conjunto \alert{pistão-cilindro};
                \item<6->  Mecanismo \alert{biela-manivela}.
            \end{itemize}
            \item<7->  Inclui parâmetros \alert{operacionais} do \alert{motor}:
            \begin{itemize}
                \item<8->  \alert{Velocidade angular} (rotação);
                \item<9->  Ângulo de \alert{ignição} e
                \item<10-> \alert{Duração da combustão}.
            \end{itemize}
        \end{itemize}
    \end{frame}
    %-------------------------------------------------------------------------------------------

%-----------------------------------------------------------------------------------------------
\section{Tópicos de Leitura}
%-----------------------------------------------------------------------------------------------

    %------------------------------------------------------------------------------------------
    \begin{frame}[allowframebreaks]{Tópicos de Leitura}
        \begin{thebibliography}{Brunetti, F., 2012}
            \bibitem[Brunetti, F., 2012]{2012-BrunettiF-Blucher}
                Brunetti, F.
                \newblock{%
                    {\em Motores de combustão interna\/}.
                    \alert{Capítulos 1 e 2.}
                }
                \newblock{\footnotesize Blücher. São Paulo. ISBN 978-85-2120-708-5.}
            \bibitem[Silva, R.~K.~de O., 2018]{2018-SilvaRKO-UTFPR}
                Silva, R.~K.~de O.
                \newblock{%
                    {\em Modelo ar-combustível de tempo finito de adição de calor de motores
                    {O}tto\/}.
                }
                \newblock{\alert{Repositório Roca UTFPR.}}
                \newblock{\footnotesize repositorio.roca.utfpr.edu.br/jspui/handle/1/8786.}
        \end{thebibliography}
    \end{frame}
    %------------------------------------------------------------------------------------------

    % Finishes with stunning image, with credit
    \ArtEndW{pexels-kaboompics-com-5969}{jpg}{txt}

%-----------------------------------------------------------------------------------------------
\end{document}
%-----------------------------------------------------------------------------------------------

